\documentclass[../../../../SoftwareRequirementsSpecification.tex]{subfiles}
\begin{document}
% Richiede le macro definite in src/macros/useCaseMacros.tex
% (incluse nel preambolo di SoftwareRequirementsSpecification.tex)

\ucstart
\begin{tabularx}{\textwidth}{|l|X|}
  \hline
  \ucid{A-02} \\ \hline
  \ucname{Registrazione Allenamento Svolto} \\ \hline
  \ucactor{Atleta} \\ \hline
  \ucdescription{L'atleta registra ciò che ha
    effettivamente svolto in una sessione di allenamento della
    scheda che sta seguendo, confermando o correggendo i valori che
    gli erano stati prescritti.} \\ \hline
  \ucprecond{L'atleta deve essere autenticato.} \\ \hline
  \ucpostcond{L'allenamento svolto risulta
    registrato ed è consultabile dall'atleta (caso d'uso A-04) e dal
    proprio PT (caso d'uso PT-08).} \\ \hline
  \ucmainflow{
    \item L'atleta atterra sulla pagina di registrazione di un
      allenamento.
    \item Il sistema propone la data odierna.
    \item L'atleta conferma la data proposta oppure ne indica
      un'altra (vedi Nota 1).
    \item Il sistema individua la scheda assegnata che copre quella
      data, ne ricava la settimana di allenamento e mostra le
      sessioni previste in quella settimana, segnalando quelle per
      cui l'atleta ha già registrato un allenamento.
    \item L'atleta preme sulla sessione che ha svolto.
    \item Il sistema mostra l'elenco degli esercizi prescritti per
      quella sessione in quella settimana, riportando per ciascuno
      (vedi Note 2 e 3):
      \begin{itemize}
        \item il peso sollevato;
        \item il numero di serie e di ripetizioni eseguite;
        \item la possibilità di marcare l'esercizio come saltato.
      \end{itemize}
    \item L'atleta conferma o corregge i valori di ciascun esercizio
      e marca come saltati quelli che non ha svolto.
    \item L'atleta preme il tasto "Salva".
    \item Il sistema registra l'allenamento svolto (vedi Nota 7).
  } \\ \hline
  \ucaltflow{
    \textbf{3a. Data futura} \newline
    - L'atleta indica una data successiva a quella odierna. \newline
    - Il sistema mostra un messaggio ("Non è possibile registrare un
    allenamento per una data futura"). \newline
    - Il flusso riprende dal punto 3. \newline
    \textbf{3b. Nessuna scheda assegnata in quella data} \newline
    - Nessuna delle schede assegnate all'atleta copre la data
    indicata. \newline
    - Il sistema mostra un messaggio ("Nessuna scheda assegnata per
    la data indicata"). \newline
    - Il flusso riprende dal punto 3. \newline
    \textbf{3c. Nessuna scheda assegnata} \newline
    - All'atleta non è mai stata assegnata alcuna scheda. \newline
    - Il sistema mostra un messaggio ("Nessuna scheda di allenamento
    assegnata"). \newline
    - Il flusso termina. \newline
    \textbf{5a. Allenamento già registrato} \newline
    - L'atleta preme su una sessione per cui ha già registrato un
    allenamento in quella settimana (vedi Nota 6). \newline
    - Il sistema mostra l'elenco degli esercizi precompilato con i
    valori che l'atleta aveva registrato, anziché con quelli
    prescritti; il salvataggio sostituisce l'allenamento
    registrato. \newline
    - Il flusso riprende dal punto 7. \newline
    \textbf{5b. Cancellazione di un allenamento registrato} \newline
    - Nella situazione del flusso 5a l'atleta preme il tasto
    "Elimina" anziché correggere i valori. \newline
    - Il sistema chiede conferma e, ottenutala, cancella
    l'allenamento registrato. \newline
    - Il flusso termina. \newline
    \textbf{8a. Valore obbligatorio non indicato} \newline
    - Su un esercizio non marcato come saltato manca un valore la cui
    compilazione è obbligatoria (vedi Nota 3). \newline
    - Il sistema mostra un messaggio che segnala gli esercizi
    incompleti e non registra nulla. \newline
    - Il flusso riprende dal punto 7.
  } \\ \hline
  \ucnote{
    \textbf{Nota 1:} la data indicata non può essere successiva a
    quella odierna e deve cadere nel periodo di validità di una
    scheda assegnata all'atleta. Registrare in ritardo è ammesso,
    anche dopo che la scheda è terminata, purché la data
    dell'allenamento cada in quel periodo. La settimana di
    allenamento non è chiesta all'atleta: il sistema la ricava dalla
    data e dalla data di inizio dell'assegnazione. Poiché i periodi
    di validità delle schede assegnate a uno stesso atleta non si
    sovrappongono, una data individua al più una scheda
    assegnata.\newline
    \textbf{Nota 2:} gli esercizi mostrati al passo 6 e i valori con
    cui sono precompilati sono quelli effettivamente prescritti a
    quell'atleta per quella settimana, cioè quelli che il caso d'uso
    A-04 gli mostra: derivano dagli allenamenti configurati nel
    blocco mensile (PT-04), dalla progressione di ciascun esercizio,
    dai pesi inseriti dal PT al momento dell'assegnazione (PT-05) e
    dalle personalizzazioni eventualmente introdotte dal PT (PT-09).
    Un superset è mostrato come i singoli esercizi che lo
    compongono.\newline
    \textbf{Nota 3:} i valori registrati sono sempre numeri singoli.
    Dove la scheda prescrive un intervallo di ripetizioni o
    l'esecuzione a cedimento non c'è alcun numero da precompilare: il
    campo è mostrato vuoto e la sua compilazione è obbligatoria,
    perché è l'unico dato che il sistema non può ricavare da sé. Il
    peso non è richiesto per gli esercizi a corpo libero. Non è
    previsto registrare il tempo di recupero effettivamente
    osservato.\newline
    \textbf{Nota 4:} l'atleta non può dichiarare di avere sostituito
    un esercizio con un altro: può soltanto marcarlo come saltato.
    La sostituzione di un esercizio è una modifica di ciò che è
    prescritto e spetta al solo PT (flusso 8b del caso d'uso
    PT-09).\newline
    \textbf{Nota 5:} l'assenza di un allenamento registrato significa
    che l'atleta non lo ha registrato, non che non lo ha svolto: il
    sistema non può distinguere i due casi. Una sessione i cui
    esercizi siano tutti marcati come saltati è invece una sessione
    che l'atleta non ha svolto.\newline
    \textbf{Nota 6:} per ciascuna sessione di ciascuna settimana di
    una scheda assegnata esiste al più un allenamento registrato: la
    scheda prescrive ogni sessione una volta sola all'interno della
    settimana.\newline
    \textbf{Nota 7:} l'allenamento registrato può essere corretto o
    cancellato dall'atleta senza limiti di tempo. Il PT non può
    crearlo, modificarlo né cancellarlo: si limita a consultarlo
    (caso d'uso PT-08). Il salvataggio non invia alcuna notifica al
    PT, che consulta gli allenamenti registrati quando lo ritiene
    opportuno.\newline
    \textbf{Nota 8:} il PT può personalizzare la settimana corrente
    (caso d'uso PT-09) anche dopo che l'atleta vi ha registrato un
    allenamento. L'allenamento registrato non cambia: cambia la
    prescrizione con cui viene confrontato. Il caso non si presenta
    per le settimane già trascorse, che PT-09 non rende
    personalizzabili.
  } \\ \hline
\end{tabularx}
\end{document}
